%=============================================================
% Calibration Degradation in Fused ECFP Pipelines
% and Temperature Scaling Mitigation
% Technical Report — Buddharak Sodtana
%=============================================================
\documentclass[11pt,a4paper,twocolumn]{article}

\usepackage[top=0.85in, bottom=0.85in, left=0.7in, right=0.7in]{geometry}
\setlength{\columnsep}{18pt}

\usepackage[T1]{fontenc}
\usepackage[utf8]{inputenc}
\usepackage{setspace}
\usepackage{amsmath,amssymb}
\usepackage{booktabs}
\usepackage{array}
\usepackage{tabularx}
\usepackage{graphicx}
\usepackage{hyperref}
\usepackage{float}
\usepackage{enumitem}
\usepackage{caption}
\usepackage{dblfloatfix}
\usepackage{listings}
\usepackage{xcolor}
\usepackage{needspace}

\hypersetup{
  colorlinks=true,
  linkcolor=blue,
  urlcolor=blue,
  pdftitle={Calibration Degradation in Fused ECFP Pipelines},
  pdfauthor={Buddharak Sodtana}
}

\lstset{
  basicstyle=\ttfamily\scriptsize,
  keywordstyle=\color{blue!70!black}\bfseries,
  commentstyle=\color{green!50!black}\itshape,
  stringstyle=\color{red!60!black},
  backgroundcolor=\color{gray!8},
  frame=single,
  framerule=0.4pt,
  rulecolor=\color{gray!40},
  breaklines=true,
  breakatwhitespace=false,
  showstringspaces=false,
  tabsize=4,
  captionpos=b,
  aboveskip=6pt,
  belowskip=4pt,
  xleftmargin=2pt,
  xrightmargin=2pt,
  columns=fullflexible
}

\newcolumntype{Y}{>{\raggedright\arraybackslash}X}
\renewcommand{\arraystretch}{1.10}
\setstretch{1.05}
\setlist{nosep,leftmargin=*}

\title{Calibration Degradation in Fused ECFP Pipelines\\
       and Temperature Scaling Mitigation}
\author{Buddharak Sodtana\\[0.5em]
        \small Technical Report for Cold-Drug DDI Benchmark}
\date{March 2026}

\begin{document}

%-------------------------------------------------------------
% FRONT MATTER — single column
%-------------------------------------------------------------
\onecolumn
\maketitle

\begin{abstract}
Calibration degradation after fusing Extended-Connectivity Fingerprint
(ECFP) features is a predictable side-effect of increased effective model
capacity, heterogeneous feature sources that create conflicting evidence and
region-dependent confidence errors, and distributional shifts that the
fusion model amplifies rather than averages out.
This report provides a comprehensive analysis of why calibration worsens
when ECFP4 fingerprints are fused with graph attention network embeddings
in the context of cold-drug DDI prediction, and presents a systematic
mitigation strategy centered on post-hoc temperature scaling.
We define key concepts, identify six mechanistic causes of calibration
failure, survey the relevant evidence base, recommend diagnostic metrics
and chart types, and provide an implementation blueprint for temperature
scaling on fused ECFP outputs.
The report concludes with practical recommendations for writing a
calibration-aware evaluation chapter.

\noindent\textbf{Keywords:} calibration, ECFP, temperature scaling,
drug--drug interaction, multimodal fusion, reliability diagram
\end{abstract}

\tableofcontents
\vspace{8pt}
\twocolumn

%=============================================================
\section{Executive Summary}
%=============================================================

A tight, defensible mitigation strategy involves three stages.
\textit{First}, diagnose \textit{how} calibration breaks: global
overconfidence vs.\ classwise vs.\ subsetwise; i.i.d.\ vs.\ shifted.
Use reliability diagrams plus scalar metrics (ECE, MCE, Brier, NLL) and
report binning choices and uncertainty around estimates.
\textit{Second}, apply post-hoc temperature scaling to the final fused
logits first; it is low-variance (1 parameter), preserves accuracy and
argmax in multiclass settings, and is a standard baseline in modern
calibration work.
\textit{Third}, if miscalibration is not mostly global---e.g., it varies
by class, interaction type, cold-start subset, or confidence
level---use structured extensions: classwise or vector scaling, adaptive
or local temperature scaling, or (when calibration data suffice) isotonic
or Dirichlet-style calibrators.
Calibration must be validated under the fusion regimes of interest (e.g.,
cold-drug splits, temporally shifted splits, scaffold splits), because
calibration that looks good i.i.d.\ can fail under shift; ensembles are
often more robust than post-hoc scaling under shift, but at higher
compute cost.

%=============================================================
\section{Definitions and Scope}
%=============================================================

\subsection{Extended-Connectivity Fingerprints (ECFP)}
ECFP are circular, neighbourhood-expansion fingerprints designed for
structure--activity modelling.
They iterate over atom environments up to a specified radius, hash
identifiers for local substructures, and fold them into fixed-length sparse
bit or count vectors (e.g., 1024/2048/4096 bits).
A key implementation detail for the fusion--calibration relationship:
\textbf{folding increases hash collisions}, i.e., distinct substructures
mapping to the same bit position, introducing structured noise and
sometimes spurious correlations that a high-capacity fusion model can
exploit with high confidence.

\subsection{Unbindable Data Types}
In this report, \textbf{unbindable data types} refers to two feature sources
that cannot be meaningfully bound by a fixed, semantics-preserving
alignment without learning a mapping or choosing an arbitrary binding
operator.
This includes two common chemical ML cases:
\begin{itemize}
\item \textbf{Pairwise entity fusion:} two molecules (drug A, drug B) each
      encoded as ECFP vectors, where the pair is unordered but representing
      it requires an ordering or a permutation-invariant operator.
\item \textbf{Heterogeneous fusion:} ECFP (hashed sparse bits/counts)
      fused with learned graph embeddings, sequence embeddings, or
      knowledge-graph features, where feature scales, sparsity, and
      semantics differ.
\end{itemize}
This aligns with multimodal ML's core challenge: heterogeneity across
modalities implies non-trivial alignment and fusion choices; naive fusion
can create brittle decision rules and confidence behaviour.

\subsection{Calibration}
A probabilistic classifier is \textbf{calibrated} if predicted
probabilities match empirical correctness frequencies (e.g., among
predictions with confidence 0.8, approximately 80\% are correct).
Reliability diagrams visualise deviations; proper scoring rules (log loss,
Brier) quantify probabilistic accuracy.
In modern deep models, high accuracy does not imply calibration;
architecture and training choices can systematically produce overconfident
probabilities.

\subsection{Temperature Scaling}
Temperature scaling is a post-hoc method that rescales logits $\mathbf{z}$
by a single positive scalar $T$:
\begin{equation}
  p_T(y\!=\!c \mid x) = \mathrm{softmax}\!\left(\frac{z_c(x)}{T}\right).
\end{equation}
$T$ is fitted on a held-out calibration set by minimising negative
log-likelihood (NLL).
In multiclass classification, it preserves argmax predictions (hence
usually preserves accuracy) while adjusting confidence.
It is a single-parameter special case of scaling approaches related to
Platt-style calibration.

%=============================================================
\section{Mechanisms of Calibration Degradation}
\label{sec:mechanisms}
%=============================================================

Calibration degradation in fusion is rarely mysterious; it typically falls
out of the interaction between capacity, evidence aggregation, and
mismatch/noise.

\subsection{Statistical and Optimisation Causes}
Cross-entropy training with hard labels pushes models to create large logit
margins, yielding very confident softmax outputs even when the underlying
problem contains ambiguity (aleatoric uncertainty) or the model is
misspecified.
Increasing effective capacity (depth/width) and certain training components
(e.g., BatchNorm) have been empirically associated with worse calibration
despite strong accuracy.
Fusion often increases effective capacity (more inputs, more parameters,
richer interactions), so it can improve discrimination while producing
\textit{sharper} probabilities---often overconfident unless regularised or
calibrated.

\subsection{Representational Causes Specific to ECFP}
ECFP's hashed/folded representation introduces two properties that matter
for calibration when fused:
\begin{itemize}
\item \textbf{Collision-induced ambiguity:} folding raises collision
      probability, causing multiple distinct substructures to share bits,
      which is a form of structured feature noise.
\item \textbf{Sparsity and heavy-tailed feature counts:} bit vectors
      create sparse, high-variance activation patterns.
      When a fusion model learns that certain sparse patterns correlate
      with labels, it can become extremely confident on those patterns,
      even if the correlation is partly artefactual.
\end{itemize}
The net effect is that adding ECFP to another representation can increase
the \textit{perceived separability} of the training data, which increases
confidence more than it increases true conditional correctness.

\subsection{Feature Overlap and Evidence Double Counting}
When two sources encode overlapping information (e.g., graph features and
ECFP both derived from 2D structure), the model may \textbf{double count
correlated evidence}, inflating confidence.
Conversely, when sources conflict (e.g., noisy assay features vs.\
structure), naive fusion can still output high confidence because the model
learns spurious rules for resolving conflicts.
Fusion increases representation richness but also increases the space of
shortcuts that yield large margins.

\subsection{Distributional Shift and Split-Induced Mismatch}
Calibration is distribution-dependent.
Even if you temperature-scale on an i.i.d.\ validation split, calibration
can deteriorate under chemical novelty (scaffold/cold-start splits),
temporal drift (newer chemistry), or domain shift between datasets.
This is particularly salient in DDI settings where new drugs differ from
known drugs.
Large-scale evaluations of predictive uncertainty under dataset shift show
that post-hoc calibration can fall short under shift, and methods that
marginalise over models (e.g., deep ensembles) are often more robust.

\subsection{Label Noise and Incomplete Labels}
Cheminformatics and pharmacology labels commonly contain measurement noise,
extraction/curation noise, or incompleteness (unknown interactions treated
as negatives).
Hard-label cross-entropy can translate such noise into overconfident
predictions; empirical work on mixup/label smoothing argues that
overconfidence is partly a consequence of hard targets.

%=============================================================
\section{Summary of Causes, Signatures, and Mitigations}
%=============================================================

Table~\ref{tab:causes} summarises the six main causes of calibration
degradation, their diagnostic signatures, and recommended mitigations
ordered from simplest to most complex.

\begin{table*}[tp]
\centering
\caption{Causes of calibration degradation in fused ECFP pipelines.}
\label{tab:causes}
\resizebox{\textwidth}{!}{%
\begin{tabular}{p{3.5cm}p{4cm}p{4cm}p{5.5cm}}
\toprule
\textbf{Cause} & \textbf{Calibration Signature} &
\textbf{Diagnostic Test} & \textbf{Mitigation (simplest first)} \\
\midrule
Increased capacity via fusion &
  Global overconfidence; reliability curve below diagonal &
  Compare max-softmax histograms; NLL vs accuracy gap &
  Temperature scaling; stronger regularisation; label smoothing/mixup;
  ensembles \\
\addlinespace
ECFP folding / collisions &
  Overconfident spikes on certain sparse patterns &
  Compare calibration by fingerprint density / collision proxy;
  ablate fpSize &
  Increase fpSize; count fingerprints; reduce folding; post-hoc
  calibration \\
\addlinespace
Feature overlap (double counting) &
  Very high confidence with modest accuracy gains &
  Calibrate separately per modality branch; measure modality agreement &
  Late fusion + (re)calibration; gating with regularisation; ensemble
  fusion \\
\addlinespace
Modality conflict / unbindable fusion &
  Region-dependent miscalibration; classwise errors &
  Conditional ECE by subset (cold-start vs warm, scaffold bins) &
  Classwise/vector scaling; local/adaptive TS; isotonic/Dirichlet
  calibration \\
\addlinespace
Distribution shift &
  Calibration good i.i.d., bad OOD &
  Shifted test sets (temporal/scaffold/distribution change) &
  Ensembles; recalibration per domain; conformal prediction as backup \\
\addlinespace
Label noise / incompleteness &
  Persistent miscalibration even after TS &
  Inspect label reliability; compare soft-label training &
  Label smoothing/mixup; Bayesian/noise-aware modelling; careful
  negatives \\
\bottomrule
\end{tabular}%
}
\end{table*}

%=============================================================
\section{Evidence Base from Literature}
%=============================================================

\subsection{Calibration in Neural Networks}
Large empirical studies establish that many modern neural networks are
miscalibrated and that temperature scaling is a strong, simple post-hoc
baseline.
Guo et al.\ demonstrated that as neural network depth and width increase,
calibration degrades despite accuracy improvements, and that a single
learned temperature parameter can substantially reduce ECE.
Subsequent work revisiting calibration continues to treat temperature
scaling as a key baseline across model families.

\subsection{Multimodal and Fusion-Specific Evidence}
Fusion can worsen trustworthiness properties even when it improves
conventional performance.
Dedicated multimodal calibration work formalises desiderata such as
``confidence should not increase when one modality is removed'' and
evaluates calibration failures in multimodal settings.
Recent clinical multimodal analyses similarly find that multimodal fusion
can degrade selective prediction reliability.

\subsection{Cheminformatics: Uncertainty and Calibration}
In drug discovery contexts, multiple studies emphasise that uncertainty
quality (including calibration) matters because experimental follow-up is
expensive and models can be overconfident on out-of-distribution chemistry.
A comprehensive calibration study for neural network structure--activity
models reports that model selection strategy and calibration approaches
significantly affect reliability of uncertainty estimates.
Ensemble-based uncertainty and recalibration frameworks in molecular
property prediction show that recalibration can improve uncertainty quality
on held-out data, including under distribution shift, but at higher
training cost.

\subsection{DDI Benchmarks as a Concrete Fusion Setting}
DDI prediction is a canonical scenario where representations of two
entities (drug A, drug B) and often multiple feature types
(structure/fingerprints/targets/enzymes) are fused.
Benchmarking under distribution changes between known and new drug sets
emphasises that evaluation must consider shift.
Practical fingerprint-only DDI work explicitly tests pairwise fusion
operators (addition, subtraction, Hadamard product) for representing drug
pairs, directly relevant to unbindable entity fusion of two ECFP vectors.

%=============================================================
\section{Calibration Diagnostics and Metrics}
%=============================================================

\subsection{Reliability Diagrams}
A reliability diagram bins predictions by confidence and plots empirical
accuracy vs.\ predicted confidence.
The diagonal indicates perfect calibration; below-diagonal indicates
overconfidence; above-diagonal indicates underconfidence.

\subsection{Scalar Calibration Metrics}

Table~\ref{tab:metrics} summarises the scalar metrics with practical
guidance for reporting.

\begin{table*}[tp]
\centering
\caption{Calibration metrics: properties and reporting guidance.}
\label{tab:metrics}
\resizebox{\textwidth}{!}{%
\begin{tabular}{p{2cm}p{3cm}p{3cm}p{3.5cm}p{4.5cm}}
\toprule
\textbf{Metric} & \textbf{What It Captures} & \textbf{Strengths} &
\textbf{Caveats} & \textbf{Recommended Reporting} \\
\midrule
Reliability diagram &
  Shape of miscalibration vs confidence &
  Diagnoses over/underconfidence; communicable &
  Bin choice; hides sample counts &
  Plot with bin counts overlay; stratify by subset (OOD/cold-start) \\
\addlinespace
ECE &
  Average calibration gap &
  Simple scalar; widely used &
  Binning bias; can mask local errors &
  Specify binning (equal-width vs equal-mass), \#bins, CI via bootstrap \\
\addlinespace
MCE &
  Worst-bin gap &
  Sensitive to failure regions &
  High variance with small bins &
  Report with ECE; show where worst bin occurs \\
\addlinespace
Brier &
  Probabilistic accuracy (squared error) &
  Proper score; decomposable &
  Conflates calibration + refinement &
  Report macro and micro variants; alongside NLL \\
\addlinespace
NLL &
  Proper score aligned with TS fitting &
  Directly optimised in TS &
  Sensitive to label noise &
  Report pre/post TS; stratify by subset; include ECE diagrams \\
\bottomrule
\end{tabular}%
}
\end{table*}

\subsection{Recommended Chart Set}
For a calibration-focused evaluation, a defensible minimal set is:
\begin{enumerate}
\item Reliability diagram (before/after calibration),
\item Bar chart of ECE and MCE across methods,
\item Line or dot chart of Brier and NLL across methods,
\item Stratified ECE (e.g., by scaffold bins, cold-start tiers, or time
      bins).
\end{enumerate}

%=============================================================
\section{Temperature Scaling and Alternatives}
\label{sec:ts}
%=============================================================

\subsection{Temperature Scaling as the Required Baseline}
Temperature scaling is well suited when fusion produces mostly global
overconfidence---a common pattern when adding features increases margin
without proportionally increasing correctness.
Key properties:
\begin{itemize}
\item \textbf{Low variance / low overfitting risk}: one parameter $T$,
      fitted on a held-out calibration set.
\item \textbf{Argmax-preserving for multiclass}: rankings of logits are
      unchanged by division by $T$, so top-1 accuracy typically stays
      constant.
\item \textbf{Optimises NLL}: connects cleanly to proper scoring rules.
\end{itemize}

\subsection{Variants and Alternatives}

Table~\ref{tab:calibrators} contrasts temperature scaling variants and
common alternative calibrators.

\begin{table*}[tp]
\centering
\caption{Calibration methods: comparison for fused ECFP pipelines.}
\label{tab:calibrators}
\resizebox{\textwidth}{!}{%
\begin{tabular}{p{3.2cm}p{1.5cm}p{1cm}p{2cm}p{2cm}p{4.5cm}}
\toprule
\textbf{Method} & \textbf{Space} & \textbf{Params} &
\textbf{Argmax?} & \textbf{Data Need} & \textbf{Best For} \\
\midrule
Temperature scaling (TS) & logits & 1 & Yes & Small &
  Global over/underconfidence \\
Vector / classwise scaling & logits & $O(C)$ & Usually & Moderate &
  Class-dependent bias \\
Matrix scaling & logits & $O(C^2)$ & Not guaranteed & Large &
  Rich miscalibration (risky with few classes) \\
Local / adaptive TS & logits & $>$1 & Often & Moderate--large &
  Region-dependent miscalibration \\
Platt scaling & score$\to$prob & 2 & N/A & Small--moderate &
  Binary calibration; margin-based models \\
Isotonic regression & score$\to$prob & nonparam. & N/A & Large &
  Complex monotone calibration curves \\
Dirichlet calibration & probs/logits & $O(C^2)$ & Not guaranteed &
  Moderate--large & Multiclass calibration beyond single $T$ \\
Deep ensembles & predictive dist. & $K\times$model & N/A & Large compute &
  Robustness under dataset shift \\
\bottomrule
\end{tabular}%
}
\end{table*}

\subsection{Fusion-Aware Calibration Strategies}
For fusing two unbindable sources, calibration often improves if you
calibrate in stages:
\begin{enumerate}
\item \textbf{Branch calibration}: calibrate per-modality logits (or
      per-entity outputs) before fusion.
\item \textbf{Fusion calibration}: fuse (e.g., logit averaging or learned
      weights), then apply final TS.
\end{enumerate}
A practical variant for DDI/pair tasks with two molecules: use
permutation-invariant pair representations (e.g.,
$\mathbf{f}_a+\mathbf{f}_b$, $|\mathbf{f}_a-\mathbf{f}_b|$,
$\mathbf{f}_a\odot\mathbf{f}_b$) to reduce spurious ordering effects,
then calibrate.

%=============================================================
\section{Experimental Design Blueprint}
\label{sec:experiment}
%=============================================================

This section gives a deployment-ready experimental plan that directly
targets calibration degradation from fused ECFP outputs and tests
temperature scaling as the first mitigation.

\subsection{Candidate Datasets and Splits}

Table~\ref{tab:datasets} proposes datasets covering the regimes where
calibration often fails.

\begin{table*}[tp]
\centering
\caption{Recommended datasets and split strategies for calibration
         evaluation.}
\label{tab:datasets}
\resizebox{\textwidth}{!}{%
\begin{tabular}{p{3.5cm}p{2.5cm}p{5cm}p{4.5cm}}
\toprule
\textbf{Dataset} & \textbf{Task Type} &
\textbf{Why Useful for Calibration Study} &
\textbf{Recommended Splits} \\
\midrule
DrugBank multi-type DDI &
  Multiclass (86 types) &
  Natural two-entity fusion (drug A + drug B) with ECFP; many labels;
  real shift concerns &
  Drug-disjoint cold-start; distribution-change splits \\
\addlinespace
MoleculeNet (BBBP, BACE, HIV, ClinTox) &
  Binary/multitask &
  Standard chemical benchmarks; supports fingerprint baselines and
  scaffold splits &
  Scaffold split; random split for contrast \\
\addlinespace
TDC (ADMET, toxicity) &
  Classification/regression &
  Curated tasks with explicit emphasis on distribution shifts &
  Scaffold / temporal / TDC split functions \\
\addlinespace
ChEMBL target-wise activity &
  Binary per target &
  Large-scale structure--activity data; suited to time-based drift &
  Time split by assay; scaffold split \\
\bottomrule
\end{tabular}%
}
\end{table*}

\subsection{Fusion Strategies to Compare}

\textbf{For two-molecule ECFP fusion (DDI-style):}
\begin{itemize}
\item Concatenation $[\mathbf{f}_a \| \mathbf{f}_b]$ plus order
      augmentation.
\item Permutation-invariant operators: $\mathbf{f}_a+\mathbf{f}_b$,
      $|\mathbf{f}_a-\mathbf{f}_b|$,
      $\mathbf{f}_a\odot\mathbf{f}_b$.
\item Bilinear/outer-product interactions (higher capacity; likely worse
      calibration unless regularised).
\end{itemize}

\textbf{For heterogeneous fusion (ECFP + learned embedding):}
\begin{itemize}
\item Early fusion: concat at feature level into an MLP head.
\item Late fusion: separate heads produce logits, then combine and apply
      TS post-hoc.
\item Gated fusion: learn per-sample modality weights (high risk for
      miscalibration if gate becomes overconfident).
\end{itemize}

\subsection{Recommended Experimental Matrix}

Table~\ref{tab:matrix} defines the minimal but informative experimental
design.

\begin{table*}[tp]
\centering
\caption{Experimental matrix for calibration evaluation.}
\label{tab:matrix}
\resizebox{\textwidth}{!}{%
\begin{tabular}{p{3cm}p{3cm}p{3.5cm}p{3cm}p{4cm}}
\toprule
\textbf{Component} & \textbf{Baseline} & \textbf{Fusion Variants} &
\textbf{Calibration} & \textbf{What You Learn} \\
\midrule
Representation &
  ECFP only &
  ECFP(A)+ECFP(B) or ECFP + embedding &
  None &
  Whether fusion improves discrimination \\
\addlinespace
Pair fusion operator &
  Concat &
  Sum / abs-diff / Hadamard &
  TS &
  Whether permutation invariance improves calibration--discrimination
  trade-off \\
\addlinespace
Heterogeneous fusion &
  Early concat &
  Late logit fusion; gated fusion &
  TS; adaptive/local TS &
  Whether miscalibration is global vs region-dependent \\
\addlinespace
Calibration complexity &
  TS &
  --- &
  Vector scaling; Dirichlet; isotonic &
  When single temperature is insufficient \\
\addlinespace
Shift robustness &
  i.i.d. &
  Scaffold/cold-start/time &
  TS vs ensembles &
  Whether post-hoc calibration survives shift \\
\bottomrule
\end{tabular}%
}
\end{table*}

%=============================================================
\section{Implementation Details}
\label{sec:impl}
%=============================================================

\subsection{ECFP Generation Parameters}
Recommended ranges to justify empirically:
\begin{itemize}
\item Radius: 2 (ECFP4) or 3 (ECFP6).
\item fpSize: 1024 / 2048 / 4096.
\item Include chirality: on/off depending on task.
\item Bit vs count: both worth testing if calibration is sensitive to
      sparsity.
\end{itemize}

\subsection{Temperature Scaling Optimisation}
Practical guidelines:
\begin{itemize}
\item Optimise $\log T$ to enforce $T>0$.
\item Initial $T=1$.
\item Optimiser: L-BFGS (common) or Adam.
\item Clamp or regularise $T$ to a sane range such as $[0.05, 20]$ to
      avoid numerical pathologies.
\item Use float64 during calibration fitting if logits are large.
\end{itemize}

\subsection{Pseudocode: Multiclass Temperature Scaling}

\begin{lstlisting}[language={},caption={Temperature scaling procedure.}]
Input:
  trained_model f(x) -> logits z in R^C
  calibration dataset D_cal = {(x_i, y_i)}
Goal:
  fit scalar T > 0 to minimise NLL on D_cal
Procedure:
  1) Freeze f
  2) For each x_i in D_cal:
       z_i = f(x_i)           # logits
  3) Optimise logT:
       T = exp(logT)
       p_i = softmax(z_i / T)
       loss = mean_i( -log p_i[y_i] )
  4) Return calibrated predictor:
       p_T(x) = softmax(f(x) / T)
\end{lstlisting}

\subsection{PyTorch Implementation}

\begin{lstlisting}[language=Python,caption={Minimal temperature scaling
in PyTorch.}]
import torch
import torch.nn.functional as F

class TemperatureScaler(torch.nn.Module):
    def __init__(self, init_temp=1.0):
        super().__init__()
        self.log_T = torch.nn.Parameter(
            torch.tensor([init_temp]).log())

    @property
    def T(self):
        return self.log_T.exp()

    def forward(self, logits):
        return logits / self.T

def fit_temperature(model, calib_loader,
                    device="cuda"):
    model.eval()
    logits_list, y_list = [], []
    with torch.no_grad():
        for x, y in calib_loader:
            x = x.to(device)
            y = y.to(device)
            logits_list.append(model(x))
            y_list.append(y)
    logits = torch.cat(logits_list)
    y = torch.cat(y_list)

    scaler = TemperatureScaler(1.0).to(device)
    opt = torch.optim.LBFGS(
        scaler.parameters(),
        lr=0.1, max_iter=50)

    def closure():
        opt.zero_grad()
        scaled = scaler(logits)
        loss = F.cross_entropy(scaled, y)
        loss.backward()
        return loss

    opt.step(closure)
    return scaler
\end{lstlisting}

\subsection{Multilabel Note}
If the fused ECFP model is multilabel (e.g., multiple DDI types can
co-occur), independent sigmoid outputs are typical.
In that case a single global $T$ on logits may be too coarse; consider
per-label temperature or Platt scaling, but guard against overfitting for
rare labels.

%=============================================================
\section{Expected Results and Interpretation}
%=============================================================

\subsection{What to Expect}
In many fused-ECFP settings:
\begin{itemize}
\item \textbf{Discrimination improves, calibration worsens} after fusion,
      because the fused model becomes more confident faster than it becomes
      more correct.
      This is consistent with findings relating miscalibration to model
      capacity and hard-label training.
\item \textbf{Temperature scaling reduces ECE/MCE and improves NLL},
      often substantially, with negligible effect on accuracy for multiclass
      tasks.
      If fitted $T>1$, the model was overconfident; $T<1$ indicates
      underconfidence.
\item Under \textbf{distribution shift}, improvements from TS may shrink
      or disappear; this supports including ensemble baselines or
      shift-aware evaluation.
\end{itemize}

\subsection{Limitations and Open Questions}
\begin{itemize}
\item \textbf{ECE is not a gold standard}: its estimate depends on binning
      and sample size; complementing ECE with proper scoring rules and
      reliability plots is important.
\item \textbf{Single-temperature models cannot fix structural
      miscalibration} (classwise, subsetwise, or confidence-conditional).
      If fusion introduces heterogeneous error modes, more expressive
      calibrators may be required.
\item \textbf{Calibration under long-tail labels} is particularly hard:
      calibration can look good overall while failing catastrophically on
      rare classes; classwise ECE and stratified analyses are necessary.
\item \textbf{Fusion-specific desiderata} (e.g., confidence should not
      increase when a modality is removed) are not captured by standard
      ECE/Brier alone.
\end{itemize}

%=============================================================
\section{Recommendations for Evaluation Writing}
\label{sec:recs}
%=============================================================

A strong evaluation narrative follows this structure:
\begin{enumerate}
\item Start with a crisp statement: ``Fusion improves discrimination but
      can degrade calibration; for deployment-like decisions we need
      calibrated probabilities.'' Ground it with core calibration references
      and at least one domain-specific reference.
\item Define the fusion setting precisely: pairwise fusion (two drugs) vs.\
      heterogeneous fusion (ECFP + embedding), and state why the sources
      are unbindable.
\item Present a short theory-of-failure table (Table~\ref{tab:causes}),
      then the experimental matrix.
\item Report calibration with \textbf{both plots and proper scores}
      (reliability diagrams + ECE/MCE + Brier + NLL), and be explicit
      about binning and calibration-set design.
\item Treat temperature scaling as the \textit{required baseline}; then
      motivate any more complex calibrator only if TS fails in identifiable
      ways (e.g., classwise miscalibration, shift-specific failures).
\end{enumerate}

\subsection{Reproducibility Checklist}
To make the evaluation credible, specify:
\begin{itemize}
\item Exact split definitions (random vs scaffold vs cold-start vs
      time/distribution-change).
\item Calibration set size and sampling scheme (especially in long-tail
      label regimes).
\item Bin count and binning strategy for ECE/MCE; bootstrap CIs.
\item Whether calibration is fit once per run, per fold, or per seed
      (recommended: per fold).
\end{itemize}

\subsection{Connection to the Cold-Drug DDI Benchmark}
In the context of the cold-drug DDI benchmark presented in the companion
paper, the calibration findings are concrete: adding ECFP4 fusion to the
GAT encoder improves macro PR-AUC by +6.1 pp and accuracy by +3.8 pp,
but ECE degrades from 0.1394 to 0.1966 (+41\% relative).
This pattern matches the ``increased capacity via fusion'' cause
(Section~\ref{sec:mechanisms}): the fused model learns sharper decision
boundaries from deterministic ECFP features, producing overconfident
softmax outputs.
Temperature scaling on the validation calibration split is the
recommended first step; if classwise analysis reveals that tail classes
(bottom 18 by training support) are disproportionately miscalibrated,
vector scaling or classwise temperature should be evaluated.

\end{document}
